\section{Metodologia}

\begin{frame}{Arquitetura da solução}
    \begin{figure}[!hbt]
        \setlength{\belowcaptionskip}{-0.5em}
        \centering
        \caption{Diagrama de arquitetura de ETL para arquivos EFD ICMS IPI.}
        \centering \includegraphics[width=0.8\textwidth,height=0.8\textheight]{imagens/ETL_Alto_Nivel.pdf}
        \\
        \small{Fonte: Elaborado pelo autor (2026).}
        \label{fig:arquitetura_geral}
    \end{figure}

\end{frame}

\begin{frame}{O Parser}
  \textbf{Responsável pela conversão do formato textual para uma representação intermediária:}
  \begin{itemize}
    \item Funciona por meio varredura sequencial, iterando linha a linha sobre o documento.
    \item Utiliza descida recursiva para fazer a inserção dos registros na estrutura em memória.
    \item Garante validade de cardinalidade e da quantia dos campos.
  \end{itemize}
\end{frame}

\begin{frame}{O Loader}
  \textbf{Responsável por orquestrar a inserção da representação intermediária no banco de dados:}
  \begin{itemize}
    \item \textbf{Varredura em profundidade:} A representação intermediária é varrida usando um algoritmo \textbf{DFS} recursivo.
    \item \textbf{Uso de \textit{dispatcher}:} Ao encontrar um registro, se utiliza do campo ``colunas'' dos metadados para fazer o despache dinâmico do método que extrai as informações daquela coluna.
    \item \textbf{Integridade primeiro:} Foco na consistência em detrimento da tolerância a falhas.
  \end{itemize}
\end{frame}
